\documentclass[11pt]{article}

\usepackage{tikz}
\usepackage{amsmath}
\usepackage{geometry}
\geometry{margin=2cm}

\usetikzlibrary{angles, quotes, arrows.meta, calc}

    \usepackage[breakable]{tcolorbox}
    \usepackage{parskip} % Stop auto-indenting (to mimic markdown behaviour)
    

    % Basic figure setup, for now with no caption control since it's done
    % automatically by Pandoc (which extracts ![](path) syntax from Markdown).
    \usepackage{graphicx}
    % Keep aspect ratio if custom image width or height is specified
    \setkeys{Gin}{keepaspectratio}
    % Maintain compatibility with old templates. Remove in nbconvert 6.0
    \let\Oldincludegraphics\includegraphics
    % Ensure that by default, figures have no caption (until we provide a
    % proper Figure object with a Caption API and a way to capture that
    % in the conversion process - todo).
    \usepackage{caption}
    \DeclareCaptionFormat{nocaption}{}
    \captionsetup{format=nocaption,aboveskip=0pt,belowskip=0pt}

    \usepackage{float}
    \floatplacement{figure}{H} % forces figures to be placed at the correct location
    \usepackage{xcolor} % Allow colors to be defined
    \usepackage{enumerate} % Needed for markdown enumerations to work
    \usepackage{geometry} % Used to adjust the document margins
    \usepackage{amsmath} % Equations
    \usepackage{amssymb} % Equations
    \usepackage{textcomp} % defines textquotesingle
    % Hack from http://tex.stackexchange.com/a/47451/13684:
    \AtBeginDocument{%
        \def\PYZsq{\textquotesingle}% Upright quotes in Pygmentized code
    }
    \usepackage{upquote} % Upright quotes for verbatim code
    \usepackage{eurosym} % defines \euro

    \usepackage{iftex}
    \ifPDFTeX
        \usepackage[T1]{fontenc}
        \IfFileExists{alphabeta.sty}{
              \usepackage{alphabeta}
          }{
              \usepackage[mathletters]{ucs}
              \usepackage[utf8x]{inputenc}
          }
    \else
        \usepackage{fontspec}
        \usepackage{unicode-math}
    \fi

    \usepackage{fancyvrb} % verbatim replacement that allows latex
    \usepackage{grffile} % extends the file name processing of package graphics
                         % to support a larger range
    \makeatletter % fix for old versions of grffile with XeLaTeX
    \@ifpackagelater{grffile}{2019/11/01}
    {
      % Do nothing on new versions
    }
    {
      \def\Gread@@xetex#1{%
        \IfFileExists{"\Gin@base".bb}%
        {\Gread@eps{\Gin@base.bb}}%
        {\Gread@@xetex@aux#1}%
      }
    }
    \makeatother
    \usepackage[Export]{adjustbox} % Used to constrain images to a maximum size
    \adjustboxset{max size={0.9\linewidth}{0.9\paperheight}}

    % The hyperref package gives us a pdf with properly built
    % internal navigation ('pdf bookmarks' for the table of contents,
    % internal cross-reference links, web links for URLs, etc.)
    \usepackage{hyperref}
    % The default LaTeX title has an obnoxious amount of whitespace. By default,
    % titling removes some of it. It also provides customization options.
    \usepackage{titling}
    \usepackage{longtable} % longtable support required by pandoc >1.10
    \usepackage{booktabs}  % table support for pandoc > 1.12.2
    \usepackage{array}     % table support for pandoc >= 2.11.3
    \usepackage{calc}      % table minipage width calculation for pandoc >= 2.11.1
    \usepackage[inline]{enumitem} % IRkernel/repr support (it uses the enumerate* environment)
    \usepackage[normalem]{ulem} % ulem is needed to support strikethroughs (\sout)
                                % normalem makes italics be italics, not underlines
    \usepackage{soul}      % strikethrough (\st) support for pandoc >= 3.0.0
    \usepackage{mathrsfs}
    

    
    % Colors for the hyperref package
    \definecolor{urlcolor}{rgb}{0,.145,.698}
    \definecolor{linkcolor}{rgb}{.71,0.21,0.01}
    \definecolor{citecolor}{rgb}{.12,.54,.11}

    % ANSI colors
    \definecolor{ansi-black}{HTML}{3E424D}
    \definecolor{ansi-black-intense}{HTML}{282C36}
    \definecolor{ansi-red}{HTML}{E75C58}
    \definecolor{ansi-red-intense}{HTML}{B22B31}
    \definecolor{ansi-green}{HTML}{00A250}
    \definecolor{ansi-green-intense}{HTML}{007427}
    \definecolor{ansi-yellow}{HTML}{DDB62B}
    \definecolor{ansi-yellow-intense}{HTML}{B27D12}
    \definecolor{ansi-blue}{HTML}{208FFB}
    \definecolor{ansi-blue-intense}{HTML}{0065CA}
    \definecolor{ansi-magenta}{HTML}{D160C4}
    \definecolor{ansi-magenta-intense}{HTML}{A03196}
    \definecolor{ansi-cyan}{HTML}{60C6C8}
    \definecolor{ansi-cyan-intense}{HTML}{258F8F}
    \definecolor{ansi-white}{HTML}{C5C1B4}
    \definecolor{ansi-white-intense}{HTML}{A1A6B2}
    \definecolor{ansi-default-inverse-fg}{HTML}{FFFFFF}
    \definecolor{ansi-default-inverse-bg}{HTML}{000000}

    % common color for the border for error outputs.
    \definecolor{outerrorbackground}{HTML}{FFDFDF}

    % commands and environments needed by pandoc snippets
    % extracted from the output of `pandoc -s`
    \providecommand{\tightlist}{%
      \setlength{\itemsep}{0pt}\setlength{\parskip}{0pt}}
    \DefineVerbatimEnvironment{Highlighting}{Verbatim}{commandchars=\\\{\}}
    % Add ',fontsize=\small' for more characters per line
    \newenvironment{Shaded}{}{}
    \newcommand{\KeywordTok}[1]{\textcolor[rgb]{0.00,0.44,0.13}{\textbf{{#1}}}}
    \newcommand{\DataTypeTok}[1]{\textcolor[rgb]{0.56,0.13,0.00}{{#1}}}
    \newcommand{\DecValTok}[1]{\textcolor[rgb]{0.25,0.63,0.44}{{#1}}}
    \newcommand{\BaseNTok}[1]{\textcolor[rgb]{0.25,0.63,0.44}{{#1}}}
    \newcommand{\FloatTok}[1]{\textcolor[rgb]{0.25,0.63,0.44}{{#1}}}
    \newcommand{\CharTok}[1]{\textcolor[rgb]{0.25,0.44,0.63}{{#1}}}
    \newcommand{\StringTok}[1]{\textcolor[rgb]{0.25,0.44,0.63}{{#1}}}
    \newcommand{\CommentTok}[1]{\textcolor[rgb]{0.38,0.63,0.69}{\textit{{#1}}}}
    \newcommand{\OtherTok}[1]{\textcolor[rgb]{0.00,0.44,0.13}{{#1}}}
    \newcommand{\AlertTok}[1]{\textcolor[rgb]{1.00,0.00,0.00}{\textbf{{#1}}}}
    \newcommand{\FunctionTok}[1]{\textcolor[rgb]{0.02,0.16,0.49}{{#1}}}
    \newcommand{\RegionMarkerTok}[1]{{#1}}
    \newcommand{\ErrorTok}[1]{\textcolor[rgb]{1.00,0.00,0.00}{\textbf{{#1}}}}
    \newcommand{\NormalTok}[1]{{#1}}

    % Additional commands for more recent versions of Pandoc
    \newcommand{\ConstantTok}[1]{\textcolor[rgb]{0.53,0.00,0.00}{{#1}}}
    \newcommand{\SpecialCharTok}[1]{\textcolor[rgb]{0.25,0.44,0.63}{{#1}}}
    \newcommand{\VerbatimStringTok}[1]{\textcolor[rgb]{0.25,0.44,0.63}{{#1}}}
    \newcommand{\SpecialStringTok}[1]{\textcolor[rgb]{0.73,0.40,0.53}{{#1}}}
    \newcommand{\ImportTok}[1]{{#1}}
    \newcommand{\DocumentationTok}[1]{\textcolor[rgb]{0.73,0.13,0.13}{\textit{{#1}}}}
    \newcommand{\AnnotationTok}[1]{\textcolor[rgb]{0.38,0.63,0.69}{\textbf{\textit{{#1}}}}}
    \newcommand{\CommentVarTok}[1]{\textcolor[rgb]{0.38,0.63,0.69}{\textbf{\textit{{#1}}}}}
    \newcommand{\VariableTok}[1]{\textcolor[rgb]{0.10,0.09,0.49}{{#1}}}
    \newcommand{\ControlFlowTok}[1]{\textcolor[rgb]{0.00,0.44,0.13}{\textbf{{#1}}}}
    \newcommand{\OperatorTok}[1]{\textcolor[rgb]{0.40,0.40,0.40}{{#1}}}
    \newcommand{\BuiltInTok}[1]{{#1}}
    \newcommand{\ExtensionTok}[1]{{#1}}
    \newcommand{\PreprocessorTok}[1]{\textcolor[rgb]{0.74,0.48,0.00}{{#1}}}
    \newcommand{\AttributeTok}[1]{\textcolor[rgb]{0.49,0.56,0.16}{{#1}}}
    \newcommand{\InformationTok}[1]{\textcolor[rgb]{0.38,0.63,0.69}{\textbf{\textit{{#1}}}}}
    \newcommand{\WarningTok}[1]{\textcolor[rgb]{0.38,0.63,0.69}{\textbf{\textit{{#1}}}}}


    % Define a nice break command that doesn't care if a line doesn't already
    % exist.
    \def\br{\hspace*{\fill} \\* }
    % Math Jax compatibility definitions
    \def\gt{>}
    \def\lt{<}
    \let\Oldtex\TeX
    \let\Oldlatex\LaTeX
    \renewcommand{\TeX}{\textrm{\Oldtex}}
    \renewcommand{\LaTeX}{\textrm{\Oldlatex}}
    % Document parameters
    % Document title
    
    
    \title{\textbf{\textit{Inheritance and Polymorphism in Python OOP throught a geometry example}}}

    \date{Mai 30, 2026}

    \author{Abdesselam Filali \\ {\small \textbf{infos@filali.net}} \\ {\small \textbf{https://github.com/Abdou-fi/python\_libraries}}}    

    
    
    
    
    
% Pygments definitions
\makeatletter
\def\PY@reset{\let\PY@it=\relax \let\PY@bf=\relax%
    \let\PY@ul=\relax \let\PY@tc=\relax%
    \let\PY@bc=\relax \let\PY@ff=\relax}
\def\PY@tok#1{\csname PY@tok@#1\endcsname}
\def\PY@toks#1+{\ifx\relax#1\empty\else%
    \PY@tok{#1}\expandafter\PY@toks\fi}
\def\PY@do#1{\PY@bc{\PY@tc{\PY@ul{%
    \PY@it{\PY@bf{\PY@ff{#1}}}}}}}
\def\PY#1#2{\PY@reset\PY@toks#1+\relax+\PY@do{#2}}

\@namedef{PY@tok@w}{\def\PY@tc##1{\textcolor[rgb]{0.73,0.73,0.73}{##1}}}
\@namedef{PY@tok@c}{\let\PY@it=\textit\def\PY@tc##1{\textcolor[rgb]{0.24,0.48,0.48}{##1}}}
\@namedef{PY@tok@cp}{\def\PY@tc##1{\textcolor[rgb]{0.61,0.40,0.00}{##1}}}
\@namedef{PY@tok@k}{\let\PY@bf=\textbf\def\PY@tc##1{\textcolor[rgb]{0.00,0.50,0.00}{##1}}}
\@namedef{PY@tok@kp}{\def\PY@tc##1{\textcolor[rgb]{0.00,0.50,0.00}{##1}}}
\@namedef{PY@tok@kt}{\def\PY@tc##1{\textcolor[rgb]{0.69,0.00,0.25}{##1}}}
\@namedef{PY@tok@o}{\def\PY@tc##1{\textcolor[rgb]{0.40,0.40,0.40}{##1}}}
\@namedef{PY@tok@ow}{\let\PY@bf=\textbf\def\PY@tc##1{\textcolor[rgb]{0.67,0.13,1.00}{##1}}}
\@namedef{PY@tok@nb}{\def\PY@tc##1{\textcolor[rgb]{0.00,0.50,0.00}{##1}}}
\@namedef{PY@tok@nf}{\def\PY@tc##1{\textcolor[rgb]{0.00,0.00,1.00}{##1}}}
\@namedef{PY@tok@nc}{\let\PY@bf=\textbf\def\PY@tc##1{\textcolor[rgb]{0.00,0.00,1.00}{##1}}}
\@namedef{PY@tok@nn}{\let\PY@bf=\textbf\def\PY@tc##1{\textcolor[rgb]{0.00,0.00,1.00}{##1}}}
\@namedef{PY@tok@ne}{\let\PY@bf=\textbf\def\PY@tc##1{\textcolor[rgb]{0.80,0.25,0.22}{##1}}}
\@namedef{PY@tok@nv}{\def\PY@tc##1{\textcolor[rgb]{0.10,0.09,0.49}{##1}}}
\@namedef{PY@tok@no}{\def\PY@tc##1{\textcolor[rgb]{0.53,0.00,0.00}{##1}}}
\@namedef{PY@tok@nl}{\def\PY@tc##1{\textcolor[rgb]{0.46,0.46,0.00}{##1}}}
\@namedef{PY@tok@ni}{\let\PY@bf=\textbf\def\PY@tc##1{\textcolor[rgb]{0.44,0.44,0.44}{##1}}}
\@namedef{PY@tok@na}{\def\PY@tc##1{\textcolor[rgb]{0.41,0.47,0.13}{##1}}}
\@namedef{PY@tok@nt}{\let\PY@bf=\textbf\def\PY@tc##1{\textcolor[rgb]{0.00,0.50,0.00}{##1}}}
\@namedef{PY@tok@nd}{\def\PY@tc##1{\textcolor[rgb]{0.67,0.13,1.00}{##1}}}
\@namedef{PY@tok@s}{\def\PY@tc##1{\textcolor[rgb]{0.73,0.13,0.13}{##1}}}
\@namedef{PY@tok@sd}{\let\PY@it=\textit\def\PY@tc##1{\textcolor[rgb]{0.73,0.13,0.13}{##1}}}
\@namedef{PY@tok@si}{\let\PY@bf=\textbf\def\PY@tc##1{\textcolor[rgb]{0.64,0.35,0.47}{##1}}}
\@namedef{PY@tok@se}{\let\PY@bf=\textbf\def\PY@tc##1{\textcolor[rgb]{0.67,0.36,0.12}{##1}}}
\@namedef{PY@tok@sr}{\def\PY@tc##1{\textcolor[rgb]{0.64,0.35,0.47}{##1}}}
\@namedef{PY@tok@ss}{\def\PY@tc##1{\textcolor[rgb]{0.10,0.09,0.49}{##1}}}
\@namedef{PY@tok@sx}{\def\PY@tc##1{\textcolor[rgb]{0.00,0.50,0.00}{##1}}}
\@namedef{PY@tok@m}{\def\PY@tc##1{\textcolor[rgb]{0.40,0.40,0.40}{##1}}}
\@namedef{PY@tok@gh}{\let\PY@bf=\textbf\def\PY@tc##1{\textcolor[rgb]{0.00,0.00,0.50}{##1}}}
\@namedef{PY@tok@gu}{\let\PY@bf=\textbf\def\PY@tc##1{\textcolor[rgb]{0.50,0.00,0.50}{##1}}}
\@namedef{PY@tok@gd}{\def\PY@tc##1{\textcolor[rgb]{0.63,0.00,0.00}{##1}}}
\@namedef{PY@tok@gi}{\def\PY@tc##1{\textcolor[rgb]{0.00,0.52,0.00}{##1}}}
\@namedef{PY@tok@gr}{\def\PY@tc##1{\textcolor[rgb]{0.89,0.00,0.00}{##1}}}
\@namedef{PY@tok@ge}{\let\PY@it=\textit}
\@namedef{PY@tok@gs}{\let\PY@bf=\textbf}
\@namedef{PY@tok@ges}{\let\PY@bf=\textbf\let\PY@it=\textit}
\@namedef{PY@tok@gp}{\let\PY@bf=\textbf\def\PY@tc##1{\textcolor[rgb]{0.00,0.00,0.50}{##1}}}
\@namedef{PY@tok@go}{\def\PY@tc##1{\textcolor[rgb]{0.44,0.44,0.44}{##1}}}
\@namedef{PY@tok@gt}{\def\PY@tc##1{\textcolor[rgb]{0.00,0.27,0.87}{##1}}}
\@namedef{PY@tok@err}{\def\PY@bc##1{{\setlength{\fboxsep}{\string -\fboxrule}\fcolorbox[rgb]{1.00,0.00,0.00}{1,1,1}{\strut ##1}}}}
\@namedef{PY@tok@kc}{\let\PY@bf=\textbf\def\PY@tc##1{\textcolor[rgb]{0.00,0.50,0.00}{##1}}}
\@namedef{PY@tok@kd}{\let\PY@bf=\textbf\def\PY@tc##1{\textcolor[rgb]{0.00,0.50,0.00}{##1}}}
\@namedef{PY@tok@kn}{\let\PY@bf=\textbf\def\PY@tc##1{\textcolor[rgb]{0.00,0.50,0.00}{##1}}}
\@namedef{PY@tok@kr}{\let\PY@bf=\textbf\def\PY@tc##1{\textcolor[rgb]{0.00,0.50,0.00}{##1}}}
\@namedef{PY@tok@bp}{\def\PY@tc##1{\textcolor[rgb]{0.00,0.50,0.00}{##1}}}
\@namedef{PY@tok@fm}{\def\PY@tc##1{\textcolor[rgb]{0.00,0.00,1.00}{##1}}}
\@namedef{PY@tok@vc}{\def\PY@tc##1{\textcolor[rgb]{0.10,0.09,0.49}{##1}}}
\@namedef{PY@tok@vg}{\def\PY@tc##1{\textcolor[rgb]{0.10,0.09,0.49}{##1}}}
\@namedef{PY@tok@vi}{\def\PY@tc##1{\textcolor[rgb]{0.10,0.09,0.49}{##1}}}
\@namedef{PY@tok@vm}{\def\PY@tc##1{\textcolor[rgb]{0.10,0.09,0.49}{##1}}}
\@namedef{PY@tok@sa}{\def\PY@tc##1{\textcolor[rgb]{0.73,0.13,0.13}{##1}}}
\@namedef{PY@tok@sb}{\def\PY@tc##1{\textcolor[rgb]{0.73,0.13,0.13}{##1}}}
\@namedef{PY@tok@sc}{\def\PY@tc##1{\textcolor[rgb]{0.73,0.13,0.13}{##1}}}
\@namedef{PY@tok@dl}{\def\PY@tc##1{\textcolor[rgb]{0.73,0.13,0.13}{##1}}}
\@namedef{PY@tok@s2}{\def\PY@tc##1{\textcolor[rgb]{0.73,0.13,0.13}{##1}}}
\@namedef{PY@tok@sh}{\def\PY@tc##1{\textcolor[rgb]{0.73,0.13,0.13}{##1}}}
\@namedef{PY@tok@s1}{\def\PY@tc##1{\textcolor[rgb]{0.73,0.13,0.13}{##1}}}
\@namedef{PY@tok@mb}{\def\PY@tc##1{\textcolor[rgb]{0.40,0.40,0.40}{##1}}}
\@namedef{PY@tok@mf}{\def\PY@tc##1{\textcolor[rgb]{0.40,0.40,0.40}{##1}}}
\@namedef{PY@tok@mh}{\def\PY@tc##1{\textcolor[rgb]{0.40,0.40,0.40}{##1}}}
\@namedef{PY@tok@mi}{\def\PY@tc##1{\textcolor[rgb]{0.40,0.40,0.40}{##1}}}
\@namedef{PY@tok@il}{\def\PY@tc##1{\textcolor[rgb]{0.40,0.40,0.40}{##1}}}
\@namedef{PY@tok@mo}{\def\PY@tc##1{\textcolor[rgb]{0.40,0.40,0.40}{##1}}}
\@namedef{PY@tok@ch}{\let\PY@it=\textit\def\PY@tc##1{\textcolor[rgb]{0.24,0.48,0.48}{##1}}}
\@namedef{PY@tok@cm}{\let\PY@it=\textit\def\PY@tc##1{\textcolor[rgb]{0.24,0.48,0.48}{##1}}}
\@namedef{PY@tok@cpf}{\let\PY@it=\textit\def\PY@tc##1{\textcolor[rgb]{0.24,0.48,0.48}{##1}}}
\@namedef{PY@tok@c1}{\let\PY@it=\textit\def\PY@tc##1{\textcolor[rgb]{0.24,0.48,0.48}{##1}}}
\@namedef{PY@tok@cs}{\let\PY@it=\textit\def\PY@tc##1{\textcolor[rgb]{0.24,0.48,0.48}{##1}}}

\def\PYZbs{\char`\\}
\def\PYZus{\char`\_}
\def\PYZob{\char`\{}
\def\PYZcb{\char`\}}
\def\PYZca{\char`\^}
\def\PYZam{\char`\&}
\def\PYZlt{\char`\<}
\def\PYZgt{\char`\>}
\def\PYZsh{\char`\#}
\def\PYZpc{\char`\%}
\def\PYZdl{\char`\$}
\def\PYZhy{\char`\-}
\def\PYZsq{\char`\'}
\def\PYZdq{\char`\"}
\def\PYZti{\char`\~}
% for compatibility with earlier versions
\def\PYZat{@}
\def\PYZlb{[}
\def\PYZrb{]}
\makeatother


    % For linebreaks inside Verbatim environment from package fancyvrb.
    \makeatletter
        \newbox\Wrappedcontinuationbox
        \newbox\Wrappedvisiblespacebox
        \newcommand*\Wrappedvisiblespace {\textcolor{red}{\textvisiblespace}}
        \newcommand*\Wrappedcontinuationsymbol {\textcolor{red}{\llap{\tiny$\m@th\hookrightarrow$}}}
        \newcommand*\Wrappedcontinuationindent {3ex }
        \newcommand*\Wrappedafterbreak {\kern\Wrappedcontinuationindent\copy\Wrappedcontinuationbox}
        % Take advantage of the already applied Pygments mark-up to insert
        % potential linebreaks for TeX processing.
        %        {, <, #, %, $, ' and ": go to next line.
        %        _, }, ^, &, >, - and ~: stay at end of broken line.
        % Use of \textquotesingle for straight quote.
        \newcommand*\Wrappedbreaksatspecials {%
            \def\PYGZus{\discretionary{\char`\_}{\Wrappedafterbreak}{\char`\_}}%
            \def\PYGZob{\discretionary{}{\Wrappedafterbreak\char`\{}{\char`\{}}%
            \def\PYGZcb{\discretionary{\char`\}}{\Wrappedafterbreak}{\char`\}}}%
            \def\PYGZca{\discretionary{\char`\^}{\Wrappedafterbreak}{\char`\^}}%
            \def\PYGZam{\discretionary{\char`\&}{\Wrappedafterbreak}{\char`\&}}%
            \def\PYGZlt{\discretionary{}{\Wrappedafterbreak\char`\<}{\char`\<}}%
            \def\PYGZgt{\discretionary{\char`\>}{\Wrappedafterbreak}{\char`\>}}%
            \def\PYGZsh{\discretionary{}{\Wrappedafterbreak\char`\#}{\char`\#}}%
            \def\PYGZpc{\discretionary{}{\Wrappedafterbreak\char`\%}{\char`\%}}%
            \def\PYGZdl{\discretionary{}{\Wrappedafterbreak\char`\$}{\char`\$}}%
            \def\PYGZhy{\discretionary{\char`\-}{\Wrappedafterbreak}{\char`\-}}%
            \def\PYGZsq{\discretionary{}{\Wrappedafterbreak\textquotesingle}{\textquotesingle}}%
            \def\PYGZdq{\discretionary{}{\Wrappedafterbreak\char`\"}{\char`\"}}%
            \def\PYGZti{\discretionary{\char`\~}{\Wrappedafterbreak}{\char`\~}}%
        }
        % Some characters . , ; ? ! / are not pygmentized.
        % This macro makes them "active" and they will insert potential linebreaks
        \newcommand*\Wrappedbreaksatpunct {%
            \lccode`\~`\.\lowercase{\def~}{\discretionary{\hbox{\char`\.}}{\Wrappedafterbreak}{\hbox{\char`\.}}}%
            \lccode`\~`\,\lowercase{\def~}{\discretionary{\hbox{\char`\,}}{\Wrappedafterbreak}{\hbox{\char`\,}}}%
            \lccode`\~`\;\lowercase{\def~}{\discretionary{\hbox{\char`\;}}{\Wrappedafterbreak}{\hbox{\char`\;}}}%
            \lccode`\~`\:\lowercase{\def~}{\discretionary{\hbox{\char`\:}}{\Wrappedafterbreak}{\hbox{\char`\:}}}%
            \lccode`\~`\?\lowercase{\def~}{\discretionary{\hbox{\char`\?}}{\Wrappedafterbreak}{\hbox{\char`\?}}}%
            \lccode`\~`\!\lowercase{\def~}{\discretionary{\hbox{\char`\!}}{\Wrappedafterbreak}{\hbox{\char`\!}}}%
            \lccode`\~`\/\lowercase{\def~}{\discretionary{\hbox{\char`\/}}{\Wrappedafterbreak}{\hbox{\char`\/}}}%
            \catcode`\.\active
            \catcode`\,\active
            \catcode`\;\active
            \catcode`\:\active
            \catcode`\?\active
            \catcode`\!\active
            \catcode`\/\active
            \lccode`\~`\~
        }
    \makeatother

    \let\OriginalVerbatim=\Verbatim
    \makeatletter
    \renewcommand{\Verbatim}[1][1]{%
        %\parskip\z@skip
        \sbox\Wrappedcontinuationbox {\Wrappedcontinuationsymbol}%
        \sbox\Wrappedvisiblespacebox {\FV@SetupFont\Wrappedvisiblespace}%
        \def\FancyVerbFormatLine ##1{\hsize\linewidth
            \vtop{\raggedright\hyphenpenalty\z@\exhyphenpenalty\z@
                \doublehyphendemerits\z@\finalhyphendemerits\z@
                \strut ##1\strut}%
        }%
        % If the linebreak is at a space, the latter will be displayed as visible
        % space at end of first line, and a continuation symbol starts next line.
        % Stretch/shrink are however usually zero for typewriter font.
        \def\FV@Space {%
            \nobreak\hskip\z@ plus\fontdimen3\font minus\fontdimen4\font
            \discretionary{\copy\Wrappedvisiblespacebox}{\Wrappedafterbreak}
            {\kern\fontdimen2\font}%
        }%

        % Allow breaks at special characters using \PYG... macros.
        \Wrappedbreaksatspecials
        % Breaks at punctuation characters . , ; ? ! and / need catcode=\active
        \OriginalVerbatim[#1,codes*=\Wrappedbreaksatpunct]%
    }
    \makeatother

    % Exact colors from NB
    \definecolor{incolor}{HTML}{303F9F}
    \definecolor{outcolor}{HTML}{D84315}
    \definecolor{cellborder}{HTML}{CFCFCF}
    \definecolor{cellbackground}{HTML}{F7F7F7}

    % prompt
    \makeatletter
    \newcommand{\boxspacing}{\kern\kvtcb@left@rule\kern\kvtcb@boxsep}
    \makeatother
    \newcommand{\prompt}[4]{
        {\ttfamily\llap{{\color{#2}[#3]:\hspace{3pt}#4}}\vspace{-\baselineskip}}
    }
    

    
    % Prevent overflowing lines due to hard-to-break entities
    \sloppy
    % Setup hyperref package
    \hypersetup{
      breaklinks=true,  % so long urls are correctly broken across lines
      colorlinks=true,
      urlcolor=urlcolor,
      linkcolor=linkcolor,
      citecolor=citecolor,
      }
    % Slightly bigger margins than the latex defaults
    
    \geometry{verbose,tmargin=1in,bmargin=1in,lmargin=1in,rmargin=1in}
    

    

\begin{document}
    
    \maketitle
    
    


Python is a multiparadigm programming language that supports
object-oriented programming (OOP) as well as procedural and functional
programming. OOP is a programming paradigm that uses objects and classes
to structure code in a way that is more intuitive and easier to manage.

\section{What is OOP?}\label{what-is-oop}

OOP stands for Object-Oriented Programming.

Python is an object-oriented language, allowing you to structure your
code using classes and objects for better organization and reusability.

\section{Advantages of OOP}\label{advantages-of-oop}

\begin{itemize}
\tightlist
\item
  Provides a clear structure to programs
\item
  Makes code easier to maintain, reuse, and debug
\item
  Helps keep your code DRY (Don't Repeat Yourself)
\item
  Allows you to build reusable applications with less code
\end{itemize}

Tip: The DRY principle means you should avoid writing the same code more
than once. Move repeated code into functions or classes and reuse it.

\section{What are Classes and Objects?}\label{what-are-classes-and-objects}

Classes and objects are the two core concepts in object-oriented
programming.

A class defines what an object should look like, and an object is
created based on that class. 

Otherwise : a class is a blueprint for creating objects (a particular data
structure), an object is an instance of a class.

For example:

\begin{table}[htbp]
\centering
\begin{tabular}{|l|p{6cm}|}
\hline
\textbf{Class} & \textbf{Objects} \\ \hline
Fruit & Apple, Banana, Mango \\ \hline
Car   & Volvo, Audi, Toyota \\ \hline
\end{tabular}
\end{table}


When you create an object from a class, it inherits all the variables
and functions defined inside that class.

The main concepts are :  
\begin{itemize}
\tightlist
\item
    Classes and objects 
\item
    The \textbf{init}() method 
\item
    The self parameter 
\item
    Properties and methods 
\item
    Inheritance and polymorphism 
\item
    Encapsulation and inner classes (will be treated in a further article)
\end{itemize}


\subsection{\texorpdfstring{The \textbf{init}()
Method}{The init() Method}}\label{the-init-method}

The \textbf{init}() method is a special method in Python classes, known
as a constructor. It is automatically called when a new object is
created from a class and allows you to initialize the object's
attributes.

\begin{Shaded}
\begin{Highlighting}[]
\KeywordTok{class}\NormalTok{ Person:}
    \KeywordTok{def} \FunctionTok{\_\_init\_\_}\NormalTok{(}\VariableTok{self}\NormalTok{, name, age):}
        \VariableTok{self}\NormalTok{.name }\OperatorTok{=}\NormalTok{ name}
        \VariableTok{self}\NormalTok{.age }\OperatorTok{=}\NormalTok{ age}

\NormalTok{p1 }\OperatorTok{=}\NormalTok{ Person(}\StringTok{"Amine"}\NormalTok{, }\DecValTok{36}\NormalTok{)}
\end{Highlighting}
\end{Shaded}

\subsection{The self Parameter}\label{the-self-parameter}

The self parameter is a reference to the current instance of the class
and is used to access variables that belong to the class.

It does not have to be named self , you can call it whatever you like,
but it has to be the first parameter of any function in the class:

\begin{Shaded}
\begin{Highlighting}[]
\KeywordTok{class}\NormalTok{ Person:}
    \KeywordTok{def} \FunctionTok{\_\_init\_\_}\NormalTok{(}\VariableTok{self}\NormalTok{, name, age):}
        \VariableTok{self}\NormalTok{.name }\OperatorTok{=}\NormalTok{ name}
        \VariableTok{self}\NormalTok{.age }\OperatorTok{=}\NormalTok{ age}
    \KeywordTok{def}\NormalTok{ myfunc(}\VariableTok{self}\NormalTok{):}
        \BuiltInTok{print}\NormalTok{(}\StringTok{"Hello my name is "} \OperatorTok{+} \VariableTok{self}\NormalTok{.name)}

\NormalTok{p1 }\OperatorTok{=}\NormalTok{ Person(}\StringTok{"Amine"}\NormalTok{, }\DecValTok{36}\NormalTok{)}
\NormalTok{p1.myfunc()}
\end{Highlighting}
\end{Shaded}

\subsection{Properties and Methods}\label{properties-and-methods}

A property is a special type of attribute that allows you to define a
method as if it were an attribute. This means you can access the method
like an attribute, without needing to use parentheses.

\begin{Shaded}
\begin{Highlighting}[]
\KeywordTok{class}\NormalTok{ Person:}
    \KeywordTok{def} \FunctionTok{\_\_init\_\_}\NormalTok{(}\VariableTok{self}\NormalTok{, name, age):}
        \VariableTok{self}\NormalTok{.name }\OperatorTok{=}\NormalTok{ name}
        \VariableTok{self}\NormalTok{.age }\OperatorTok{=}\NormalTok{ age}

    \KeywordTok{def}\NormalTok{ myfunc(}\VariableTok{self}\NormalTok{):}
        \ControlFlowTok{return} \StringTok{"Hello my name is "} \OperatorTok{+} \VariableTok{self}\NormalTok{.name}

\NormalTok{p1 }\OperatorTok{=}\NormalTok{ Person(}\StringTok{"Amine"}\NormalTok{, }\DecValTok{36}\NormalTok{)}
\BuiltInTok{print}\NormalTok{(p1.myfunc)}
\end{Highlighting}
\end{Shaded}

A method is a function that is defined inside a class and is used to
perform operations on the object.

\begin{Shaded}
\begin{Highlighting}[]
\KeywordTok{class}\NormalTok{ Person:}
    \KeywordTok{def} \FunctionTok{\_\_init\_\_}\NormalTok{(}\VariableTok{self}\NormalTok{, name, age):}
        \VariableTok{self}\NormalTok{.name }\OperatorTok{=}\NormalTok{ name}
        \VariableTok{self}\NormalTok{.age }\OperatorTok{=}\NormalTok{ age}

    \KeywordTok{def}\NormalTok{ myfunc(}\VariableTok{self}\NormalTok{):}
        \BuiltInTok{print}\NormalTok{(}\StringTok{"Hello my name is "} \OperatorTok{+} \VariableTok{self}\NormalTok{.name)}

\NormalTok{p1 }\OperatorTok{=}\NormalTok{ Person(}\StringTok{"Amine"}\NormalTok{, }\DecValTok{36}\NormalTok{)}
\NormalTok{p1.myfunc()}
\end{Highlighting}
\end{Shaded}

\subsection{Inheritance and
Polymorphism}\label{inheritance-and-polymorphism}

\subsubsection{Inheritance}\label{inheritance}

Inheritance, a powerful feature in object-oriented programming, is the
concept of creating a new class that is a modified version of an
existing class. The new class inherits all the attributes and methods of
the existing class, and you can add new attributes and methods or
override existing ones.

It allows allows a new class (called a child or subclass) to inherit
attributes and methods from an existing class (called a parent or
superclass). This promotes code reusability and establishes a natural
hierarchical relationship between classes.

\begin{Shaded}
\begin{Highlighting}[]
\KeywordTok{class}\NormalTok{ Person:}

    \KeywordTok{def} \FunctionTok{\_\_init\_\_}\NormalTok{(}\VariableTok{self}\NormalTok{, name, age):}
        \VariableTok{self}\NormalTok{.name }\OperatorTok{=}\NormalTok{ name}
        \VariableTok{self}\NormalTok{.age }\OperatorTok{=}\NormalTok{ age}

    \KeywordTok{def}\NormalTok{ informations(}\VariableTok{self}\NormalTok{):}
        \BuiltInTok{print}\NormalTok{(}\StringTok{"Hello my name is "} \OperatorTok{+} \VariableTok{self}\NormalTok{.name)}

\KeywordTok{class}\NormalTok{ Student(Person):}
    \KeywordTok{def} \FunctionTok{\_\_init\_\_}\NormalTok{(}\VariableTok{self}\NormalTok{, name, age, grade):}
        \BuiltInTok{super}\NormalTok{().}\FunctionTok{\_\_init\_\_}\NormalTok{(name, age)}
        \VariableTok{self}\NormalTok{.grade }\OperatorTok{=}\NormalTok{ grade}

    \KeywordTok{def}\NormalTok{ informations(}\VariableTok{self}\NormalTok{):}
        \BuiltInTok{print}\NormalTok{(}\StringTok{"Hello my name is "} \OperatorTok{+} \VariableTok{self}\NormalTok{.name }\OperatorTok{+} \StringTok{" and I am in grade "} \OperatorTok{+} \BuiltInTok{str}\NormalTok{(}\VariableTok{self}\NormalTok{.grade))}

\NormalTok{p1 }\OperatorTok{=}\NormalTok{ Student(}\StringTok{"Amine"}\NormalTok{, }\DecValTok{36}\NormalTok{, }\DecValTok{12}\NormalTok{)}
\NormalTok{p1.informations()}
\end{Highlighting}
\end{Shaded}

\subsubsection{polymorphism}\label{polymorphism}

Polymorphism allows you to use a unified interface to access different
types of objects. This means that you can use the same method name to
perform different tasks based on the object that is calling the method.

\begin{Shaded}
\begin{Highlighting}[]
\KeywordTok{class}\NormalTok{ Person:}
    \KeywordTok{def} \FunctionTok{\_\_init\_\_}\NormalTok{(}\VariableTok{self}\NormalTok{, name, age):}
        \VariableTok{self}\NormalTok{.name }\OperatorTok{=}\NormalTok{ name}
        \VariableTok{self}\NormalTok{.age }\OperatorTok{=}\NormalTok{ age}
    \KeywordTok{def}\NormalTok{ myfunc(}\VariableTok{self}\NormalTok{):}
        \BuiltInTok{print}\NormalTok{(}\StringTok{"Hello my name is "} \OperatorTok{+} \VariableTok{self}\NormalTok{.name)}

\KeywordTok{class}\NormalTok{ Student(Person):}
    \KeywordTok{def} \FunctionTok{\_\_init\_\_}\NormalTok{(}\VariableTok{self}\NormalTok{, name, age, grade):}
        \BuiltInTok{super}\NormalTok{().}\FunctionTok{\_\_init\_\_}\NormalTok{(name, age)}
        \VariableTok{self}\NormalTok{.grade }\OperatorTok{=}\NormalTok{ grade}

    \KeywordTok{def}\NormalTok{ myfunc(}\VariableTok{self}\NormalTok{):}
        \BuiltInTok{print}\NormalTok{(}\StringTok{"Hello my name is "} \OperatorTok{+} \VariableTok{self}\NormalTok{.name }\OperatorTok{+} \StringTok{" and I am in grade "} \OperatorTok{+} \BuiltInTok{str}\NormalTok{(}\VariableTok{self}\NormalTok{.grade))}

\NormalTok{p1 }\OperatorTok{=}\NormalTok{ Person(}\StringTok{"Amine"}\NormalTok{, }\DecValTok{36}\NormalTok{)}
\NormalTok{p2 }\OperatorTok{=}\NormalTok{ Student(}\StringTok{"Sarah"}\NormalTok{, }\DecValTok{12}\NormalTok{, }\DecValTok{10}\NormalTok{)}

\NormalTok{p1.myfunc()}
\NormalTok{p2.myfunc()}
\end{Highlighting}
\end{Shaded}


\section{Example of Inheritance \& Polymorphism}\label{example-of-inheritance-polymorphism}

The following representation is a parallelogram with sides \(\ell\) and \(w\) and height \(h\) (the distance between the two parallel sides of length \(\ell\)).
You find the relations between the angle $\theta$, the length and the width of the parallelogram translated into vertex coordinates.

\begin{center}
\begin{tikzpicture}[scale=1.4, >=Stealth]

  % A is at a general position (x0, y0), visually offset from the origin
  \coordinate (O)     at (0, 0);
  \coordinate (A)     at (1.2, 0.8);
  \coordinate (B)     at (4.2, 0.8);
  \coordinate (C)     at (5.7, 2.6);
  \coordinate (D)     at (2.7, 2.6);

  % Feet on x-axis
  \coordinate (Ax) at (1.2, 0);
  \coordinate (Bx) at (4.2, 0);
  \coordinate (Cx) at (5.7, 0);
  \coordinate (Dx) at (2.7, 0);

  % Feet on y-axis
  \coordinate (Ay) at (0, 0.8);
  \coordinate (Cy) at (0, 2.6);   % same height as D

  % Extension of AB beyond B (for right-side theta at B)
  \coordinate (Bext) at (5.5, 0.8);

  % Foot of perpendicular from D and C onto the horizontal through A
  \coordinate (Dfoot) at (2.7, 0.8);
  \coordinate (Cfoot) at (5.7, 0.8);

  % Draw axes
  \draw[->] (-0.3, 0) -- (6.8, 0) node[right] {$x$};
  \draw[->] (0, -0.3) -- (0, 3.4) node[above] {$y$};

  % Draw the parallelogram
  \draw[thick] (A) -- (B) -- (C) -- (D) -- cycle;

  % Dashed vertical lines (heights) from D and C down to line y=y0
  \draw[dashed] (D) -- (Dfoot);
  \draw[dashed] (C) -- (Cfoot);

  % Dotted projection lines to axes for A
  \draw[dotted] (A) -- (Ay);
  \draw[dotted] (A) -- (Ax);

  % Dotted projection lines to x-axis for B, C, D
  \draw[dotted] (B) -- (Bx);
  \draw[dotted] (C) -- (Cx);
  \draw[dotted] (D) -- (Dx);

  % Dotted projection lines to y-axis for C (and D, same height)
  \draw[dotted] (C) -- (Cy);

  % Label vertices
  \node[below left]  at (A) {$A$};
  \node[above left]  at (B) {$B$};
  \node[above right] at (C) {$C$};
  \node[above left]  at (D) {$D$};

  % Label sides
  \node[above] at ($(D)!0.5!(C)$) {$\ell$};
  \node[below] at ($(A)!0.5!(B)$) {$\ell$};
  \node[left]  at ($(A)!0.5!(D)$) {$w$};
  \node[right] at ($(B)!0.5!(C)$) {$w$};

  % Label h on dashed verticals
  \node[right] at ($(D)!0.5!(Dfoot)$) {$h$};
  \node[right] at ($(C)!0.5!(Cfoot)$) {$h$};

  % Label l' : horizontal distance from B to foot of C
  \draw[dashed] (B) -- (Cfoot);
  \node[below] at ($(B)!0.5!(Cfoot)$) {$\ell'$};

  % Angle theta at A
  \draw pic["$\theta$", draw=black, angle radius=0.55cm, angle eccentricity=1.5]
    {angle = B--A--D};

  % Angle theta at B (opens to the right, between extension of AB and BC)
  \draw pic["$\theta$", draw=black, angle radius=0.55cm, angle eccentricity=1.5]
    {angle = Bext--B--C};

  % --- X-axis tick marks and labels ---
  % x0 (A)
  \draw (1.2, 0.04) -- (1.2, -0.04);
  \node[below] at (Ax) {$x_0$};

  % x0+l (B)
  \draw (4.2, 0.04) -- (4.2, -0.04);
  \node[below] at (Bx) {$x_0{+}\ell$};

  % x0+l+wcos (C)
  \draw (5.7, 0.04) -- (5.7, -0.04);
  \node[below] at (Cx) {\small$x_0{+}\ell{+}w\cos\theta$};

  % x0+wcos (D)
  \draw (2.7, 0.04) -- (2.7, -0.04);
  \node[below] at (Dx) {$x_0{+}w\cos\theta$};

  % --- Y-axis tick marks and labels ---
  % y0 (A and B)
  \draw (0.04, 0.8) -- (-0.04, 0.8);
  \node[left] at (Ay) {$y_0$};

  % y0+wsin (C and D)
  \draw (0.04, 2.6) -- (-0.04, 2.6);
  \node[left] at (Cy) {$y_0{+}w\sin\theta$};

\end{tikzpicture}
\end{center}

\vspace{1em}

\noindent\textbf{Relations between $\ell'$, $h$ and $w$:}
\begin{align*}
  \cos\theta &= \frac{\ell'}{w} \quad\Rightarrow\quad \ell' = w\cos\theta \\[4pt]
  \sin\theta &= \frac{h}{w}    \quad\Rightarrow\quad h   = w\sin\theta
\end{align*}

\vspace{0.5em}


\noindent\textbf{Vertex coordinates:}
\begin{align*}
  A &= (x_0,\ y_0) \\
  B &= (x_0 + \ell,\ y_0) \\
  C &= (x_0 + \ell + w\cos\theta,\ y_0 + w\sin\theta) \\
  D &= (x_0 + w\cos\theta,\ y_0 + w\sin\theta)
\end{align*}


The parallelogram is completely defined by the parameters \(x_0\), \(y_0\), \(\ell\), \(w\), and \(\theta\). The coordinates of the vertices A, B, C, and D can be calculated using these parameters.

In the following example, we will create a base class called \texttt{Shape} and then create a subclass called \texttt{Parallelogram}
that inherits from \texttt{Shape}.

    \begin{tcolorbox}[breakable, size=fbox, boxrule=1pt, pad at break*=1mm,colback=cellbackground, colframe=cellborder]
\prompt{In}{incolor}{2}{\boxspacing}
\begin{Verbatim}[commandchars=\\\{\}]
\PY{k+kn}{import} \PY{n+nn}{matplotlib}\PY{n+nn}{.}\PY{n+nn}{pyplot} \PY{k}{as} \PY{n+nn}{plt} \PY{c+c1}{\PYZsh{} Import the necessary library for plotting}
\PY{k+kn}{from} \PY{n+nn}{math} \PY{k+kn}{import} \PY{n}{cos}\PY{p}{,} \PY{n}{sin}\PY{p}{,} \PY{n}{pi}   \PY{c+c1}{\PYZsh{} Import mathematical functions for cosine, sine, and pi}

\PY{k}{class} \PY{n+nc}{Shape}\PY{p}{:}
    \PY{k}{def} \PY{n+nf+fm}{\PYZus{}\PYZus{}init\PYZus{}\PYZus{}}\PY{p}{(}\PY{n+nb+bp}{self}\PY{p}{,} \PY{n}{x0}\PY{p}{:}\PY{n+nb}{float}\PY{p}{,} \PY{n}{y0}\PY{p}{:}\PY{n+nb}{float}\PY{p}{)}\PY{p}{:}
        \PY{n+nb+bp}{self}\PY{o}{.}\PY{n}{x0} \PY{o}{=} \PY{n}{x0}        \PY{c+c1}{\PYZsh{} Initialize the x\PYZhy{}coordinate of the shape\PYZsq{}s reference point}
        \PY{n+nb+bp}{self}\PY{o}{.}\PY{n}{y0} \PY{o}{=} \PY{n}{y0}        \PY{c+c1}{\PYZsh{} Initialize the y\PYZhy{}coordinate of the shape\PYZsq{}s reference point}
    \PY{k}{def} \PY{n+nf}{draw\PYZus{}shape}\PY{p}{(}\PY{n+nb+bp}{self}\PY{p}{)}\PY{p}{:}   \PY{c+c1}{\PYZsh{} Define a method to draw the shape, which will be overridden in subclasses}
        \PY{k}{pass}                \PY{c+c1}{\PYZsh{} This method will be overridden in the subclass  }

\PY{c+c1}{\PYZsh{} Define a class for parallelograms that inherits from the Shape class, }
\PY{c+c1}{\PYZsh{} since a parallelogram is a specific type of shape with additional properties and methods.}
\PY{k}{class} \PY{n+nc}{Parallelogram}\PY{p}{(}\PY{n}{Shape}\PY{p}{)}\PY{p}{:}
        
    \PY{k}{def} \PY{n+nf+fm}{\PYZus{}\PYZus{}init\PYZus{}\PYZus{}}\PY{p}{(}\PY{n+nb+bp}{self}\PY{p}{,} \PY{n}{x0}\PY{p}{:}\PY{n+nb}{float}\PY{p}{,} \PY{n}{y0}\PY{p}{:}\PY{n+nb}{float}\PY{p}{,} \PY{n}{length}\PY{p}{:}\PY{n+nb}{float}\PY{p}{,} \PY{n}{width}\PY{p}{:}\PY{n+nb}{float}\PY{p}{,} \PY{n}{angle}\PY{p}{:} \PY{n+nb}{float}\PY{p}{)}\PY{p}{:}
        \PY{n+nb}{super}\PY{p}{(}\PY{p}{)}\PY{o}{.}\PY{n+nf+fm}{\PYZus{}\PYZus{}init\PYZus{}\PYZus{}}\PY{p}{(}\PY{n}{x0}\PY{p}{,} \PY{n}{y0}\PY{p}{)}            \PY{c+c1}{\PYZsh{} Call the constructor of the parent class}
        \PY{c+c1}{\PYZsh{} Initialize the angle of the parallelogram in radians}
        \PY{n+nb+bp}{self}\PY{o}{.}\PY{n}{angle} \PY{o}{=} \PY{n}{angle}                  
        \PY{c+c1}{\PYZsh{} Initialize the length of the parallelogram}
        \PY{n+nb+bp}{self}\PY{o}{.}\PY{n}{length} \PY{o}{=} \PY{n}{length}                
        \PY{c+c1}{\PYZsh{} Initialize the width of the parallelogram}
        \PY{n+nb+bp}{self}\PY{o}{.}\PY{n}{width} \PY{o}{=} \PY{n}{width}                  
        \PY{c+c1}{\PYZsh{} Calculate the height of the parallelogram using the sine of the angle}
        \PY{n+nb+bp}{self}\PY{o}{.}\PY{n}{height}\PY{o}{=} \PY{n}{width} \PY{o}{*} \PY{n}{sin}\PY{p}{(}\PY{n}{angle}\PY{p}{)}     
        \PY{c+c1}{\PYZsh{} Calculate the area (superficy) of the parallelogram using the formula: area = base * height}
        \PY{n+nb+bp}{self}\PY{o}{.}\PY{n}{superficy} \PY{o}{=} \PY{n}{length} \PY{o}{*} \PY{n+nb+bp}{self}\PY{o}{.}\PY{n}{height}   
        \PY{c+c1}{\PYZsh{} Initialize the coordinates of point A (the reference point)}
        \PY{n+nb+bp}{self}\PY{o}{.}\PY{n}{point\PYZus{}A\PYZus{}x}\PY{p}{,} \PY{n+nb+bp}{self}\PY{o}{.}\PY{n}{point\PYZus{}A\PYZus{}y} \PY{o}{=} \PY{n}{x0}\PY{p}{,} \PY{n}{y0}      
        \PY{c+c1}{\PYZsh{} Calculate the coordinates of point B by adding the length to the x\PYZhy{}coordinate of point A}
        \PY{n+nb+bp}{self}\PY{o}{.}\PY{n}{point\PYZus{}B\PYZus{}x}\PY{p}{,} \PY{n+nb+bp}{self}\PY{o}{.}\PY{n}{point\PYZus{}B\PYZus{}y} \PY{o}{=} \PY{n}{x0} \PY{o}{+} \PY{n}{length}\PY{p}{,} \PY{n}{y0}    
        \PY{c+c1}{\PYZsh{} Calculate the coordinates of point C by adding the length and the horizontal component of }
        \PY{c+c1}{\PYZsh{} the width to the x\PYZhy{}coordinate of point A, and adding the vertical component of }
        \PY{c+c1}{\PYZsh{} the width to the y\PYZhy{}coordinate of point A}
        \PY{n+nb+bp}{self}\PY{o}{.}\PY{n}{point\PYZus{}C\PYZus{}x}\PY{p}{,} \PY{n+nb+bp}{self}\PY{o}{.}\PY{n}{point\PYZus{}C\PYZus{}y} \PY{o}{=} \PY{n}{x0} \PY{o}{+} \PY{n}{length} \PY{o}{+} \PY{n}{width} \PY{o}{*} \PY{n}{cos}\PY{p}{(}\PY{n}{angle}\PY{p}{)}\PY{p}{,} \PY{n}{y0} \PY{o}{+} \PY{n}{width} \PY{o}{*} \PY{n}{sin}\PY{p}{(}\PY{n}{angle}\PY{p}{)}  
        \PY{c+c1}{\PYZsh{} Calculate the coordinates of point D by adding the horizontal component of }
        \PY{c+c1}{\PYZsh{} the width to the x\PYZhy{}coordinate of point A, and adding the vertical component of }
        \PY{c+c1}{\PYZsh{} the width to the y\PYZhy{}coordinate of point A}
        \PY{n+nb+bp}{self}\PY{o}{.}\PY{n}{point\PYZus{}D\PYZus{}x}\PY{p}{,} \PY{n+nb+bp}{self}\PY{o}{.}\PY{n}{point\PYZus{}D\PYZus{}y} \PY{o}{=} \PY{n}{x0} \PY{o}{+} \PY{n}{width} \PY{o}{*} \PY{n}{cos}\PY{p}{(}\PY{n}{angle}\PY{p}{)}\PY{p}{,} \PY{n}{y0} \PY{o}{+} \PY{n}{width} \PY{o}{*} \PY{n}{sin}\PY{p}{(}\PY{n}{angle}\PY{p}{)}   


    \PY{c+c1}{\PYZsh{} Define a method to print the description of the parallelogram, }
    \PY{c+c1}{\PYZsh{} including its properties: length, width, and the starting angle (converted to degrees for better readability)}
    \PY{k}{def} \PY{n+nf}{description}\PY{p}{(}\PY{n+nb+bp}{self}\PY{p}{)}\PY{p}{:}
        \PY{n+nb}{print}\PY{p}{(}\PY{l+s+sa}{f}\PY{l+s+s2}{\PYZdq{}}\PY{l+s+s2}{This is a }\PY{l+s+si}{\PYZob{}}\PY{n+nb+bp}{self}\PY{o}{.}\PY{n+nv+vm}{\PYZus{}\PYZus{}class\PYZus{}\PYZus{}}\PY{o}{.}\PY{n+nv+vm}{\PYZus{}\PYZus{}name\PYZus{}\PYZus{}}\PY{o}{.}\PY{n}{lower}\PY{p}{(}\PY{p}{)}\PY{l+s+si}{\PYZcb{}}\PY{l+s+s2}{ with the following properties:}\PY{l+s+s2}{\PYZdq{}}\PY{p}{)}
        \PY{n+nb}{print}\PY{p}{(}\PY{l+s+s2}{\PYZdq{}}\PY{l+s+s2}{Length: }\PY{l+s+se}{\PYZbs{}t}\PY{l+s+s2}{\PYZdq{}}\PY{p}{,} \PY{n+nb+bp}{self}\PY{o}{.}\PY{n}{length}\PY{p}{)}        
        \PY{n+nb}{print}\PY{p}{(}\PY{l+s+s2}{\PYZdq{}}\PY{l+s+s2}{Width: }\PY{l+s+se}{\PYZbs{}t}\PY{l+s+se}{\PYZbs{}t}\PY{l+s+s2}{\PYZdq{}}\PY{p}{,} \PY{n+nb+bp}{self}\PY{o}{.}\PY{n}{width}\PY{p}{)}        
        \PY{c+c1}{\PYZsh{} Convert angle from radians to degrees for a more intuitive understanding}
        \PY{n+nb}{print}\PY{p}{(}\PY{l+s+s2}{\PYZdq{}}\PY{l+s+s2}{Angle (in degrees): }\PY{l+s+se}{\PYZbs{}t}\PY{l+s+s2}{\PYZdq{}}\PY{p}{,} \PY{n+nb+bp}{self}\PY{o}{.}\PY{n}{angle} \PY{o}{*} \PY{l+m+mi}{180} \PY{o}{/} \PY{n}{pi}\PY{p}{)}   
        
    \PY{c+c1}{\PYZsh{} Define the x\PYZhy{}coordinates of the vertices of the parallelogram in order (A, B, C, D, and back to A to close the shape)}
    \PY{k}{def} \PY{n+nf}{draw\PYZus{}shape}\PY{p}{(}\PY{n+nb+bp}{self}\PY{p}{)}\PY{p}{:}
        \PY{n}{data\PYZus{}x} \PY{o}{=} \PY{p}{[}\PY{n+nb+bp}{self}\PY{o}{.}\PY{n}{point\PYZus{}A\PYZus{}x}\PY{p}{,} 
                  \PY{n+nb+bp}{self}\PY{o}{.}\PY{n}{point\PYZus{}B\PYZus{}x}\PY{p}{,} 
                  \PY{n+nb+bp}{self}\PY{o}{.}\PY{n}{point\PYZus{}C\PYZus{}x}\PY{p}{,}  
                  \PY{n+nb+bp}{self}\PY{o}{.}\PY{n}{point\PYZus{}D\PYZus{}x}\PY{p}{,}  
                  \PY{n+nb+bp}{self}\PY{o}{.}\PY{n}{point\PYZus{}A\PYZus{}x} \PY{p}{]}   
        \PY{n}{data\PYZus{}y} \PY{o}{=} \PY{p}{[}\PY{n+nb+bp}{self}\PY{o}{.}\PY{n}{point\PYZus{}A\PYZus{}y}\PY{p}{,} 
                  \PY{n+nb+bp}{self}\PY{o}{.}\PY{n}{point\PYZus{}B\PYZus{}y}\PY{p}{,} 
                  \PY{n+nb+bp}{self}\PY{o}{.}\PY{n}{point\PYZus{}C\PYZus{}y}\PY{p}{,} 
                  \PY{n+nb+bp}{self}\PY{o}{.}\PY{n}{point\PYZus{}D\PYZus{}y}\PY{p}{,} 
                  \PY{n+nb+bp}{self}\PY{o}{.}\PY{n}{point\PYZus{}A\PYZus{}y} \PY{p}{]}
        
        \PY{n}{labels} \PY{o}{=} \PY{p}{[}\PY{l+s+s1}{\PYZsq{}}\PY{l+s+s1}{A}\PY{l+s+s1}{\PYZsq{}}\PY{p}{,} \PY{l+s+s1}{\PYZsq{}}\PY{l+s+s1}{B}\PY{l+s+s1}{\PYZsq{}}\PY{p}{,} \PY{l+s+s1}{\PYZsq{}}\PY{l+s+s1}{C}\PY{l+s+s1}{\PYZsq{}}\PY{p}{,} \PY{l+s+s1}{\PYZsq{}}\PY{l+s+s1}{D}\PY{l+s+s1}{\PYZsq{}}\PY{p}{]}  \PY{c+c1}{\PYZsh{} Define labels for the vertices of the parallelogram}
        \PY{k}{for} \PY{n}{i}\PY{p}{,} \PY{n}{label} \PY{o+ow}{in} \PY{n+nb}{enumerate}\PY{p}{(}\PY{n}{labels}\PY{p}{)}\PY{p}{:}
            \PY{n}{plt}\PY{o}{.}\PY{n}{text}\PY{p}{(}\PY{n}{data\PYZus{}x}\PY{p}{[}\PY{n}{i}\PY{p}{]}\PY{p}{,} \PY{n}{data\PYZus{}y}\PY{p}{[}\PY{n}{i}\PY{p}{]}\PY{p}{,} \PY{n}{label}\PY{p}{,} \PY{n}{fontsize}\PY{o}{=}\PY{l+m+mi}{10}\PY{p}{,} \PY{n}{ha}\PY{o}{=}\PY{l+s+s1}{\PYZsq{}}\PY{l+s+s1}{right}\PY{l+s+s1}{\PYZsq{}}\PY{p}{,} \PY{n}{va}\PY{o}{=}\PY{l+s+s1}{\PYZsq{}}\PY{l+s+s1}{bottom}\PY{l+s+s1}{\PYZsq{}}\PY{p}{)}  \PY{c+c1}{\PYZsh{} Annotate each vertex with its label}
        
        \PY{n}{plt}\PY{o}{.}\PY{n}{plot}\PY{p}{(}\PY{n}{data\PYZus{}x}\PY{p}{,} \PY{n}{data\PYZus{}y}\PY{p}{,} \PY{n}{linestyle}\PY{o}{=}\PY{l+s+s1}{\PYZsq{}}\PY{l+s+s1}{\PYZhy{}}\PY{l+s+s1}{\PYZsq{}}\PY{p}{)} \PY{c+c1}{\PYZsh{} Plot the vertices of the parallelogram by connecting points A, B, C, D, and back to A}
        \PY{n}{plt}\PY{o}{.}\PY{n}{axis}\PY{p}{(}\PY{l+s+s1}{\PYZsq{}}\PY{l+s+s1}{equal}\PY{l+s+s1}{\PYZsq{}}\PY{p}{)}       \PY{c+c1}{\PYZsh{} Set equal scaling for x and y axes to maintain the shape\PYZsq{}s proportions}
        \PY{n}{plt}\PY{o}{.}\PY{n}{title}\PY{p}{(}\PY{l+s+sa}{f}\PY{l+s+s2}{\PYZdq{}}\PY{l+s+si}{\PYZob{}}\PY{n+nb+bp}{self}\PY{o}{.}\PY{n+nv+vm}{\PYZus{}\PYZus{}class\PYZus{}\PYZus{}}\PY{o}{.}\PY{n+nv+vm}{\PYZus{}\PYZus{}name\PYZus{}\PYZus{}}\PY{o}{.}\PY{n}{lower}\PY{p}{(}\PY{p}{)}\PY{l+s+si}{\PYZcb{}}\PY{l+s+s2}{ representation}\PY{l+s+s2}{\PYZdq{}}\PY{p}{)} \PY{c+c1}{\PYZsh{} Set the title of the plot to the name of the shape}
        \PY{n}{plt}\PY{o}{.}\PY{n}{show}\PY{p}{(}\PY{p}{)}              \PY{c+c1}{\PYZsh{} Display the plot of the parallelogram}
        \PY{n}{plt}\PY{o}{.}\PY{n}{close}\PY{p}{(}\PY{p}{)}             \PY{c+c1}{\PYZsh{} Close the plot after displaying it to free up resources and avoid displaying multiple plots in interactive environments}


    \PY{c+c1}{\PYZsh{} Define a method to print the area (superficy) of the parallelogram, which is calculated in the constructor}
    \PY{k}{def} \PY{n+nf}{print\PYZus{}superficy}\PY{p}{(}\PY{n+nb+bp}{self}\PY{p}{)}\PY{p}{:}
        \PY{n+nb}{print}\PY{p}{(}\PY{l+s+sa}{f}\PY{l+s+s2}{\PYZdq{}}\PY{l+s+s2}{The superficy of the }\PY{l+s+si}{\PYZob{}}\PY{n+nb+bp}{self}\PY{o}{.}\PY{n+nv+vm}{\PYZus{}\PYZus{}class\PYZus{}\PYZus{}}\PY{o}{.}\PY{n+nv+vm}{\PYZus{}\PYZus{}name\PYZus{}\PYZus{}}\PY{o}{.}\PY{n}{lower}\PY{p}{(}\PY{p}{)}\PY{l+s+si}{\PYZcb{}}\PY{l+s+s2}{ is: }\PY{l+s+si}{\PYZob{}}\PY{n+nb+bp}{self}\PY{o}{.}\PY{n}{superficy}\PY{l+s+si}{\PYZcb{}}\PY{l+s+s2}{\PYZdq{}}\PY{p}{)}

\PY{n}{parallelogram1} \PY{o}{=} \PY{n}{Parallelogram}\PY{p}{(}\PY{l+m+mi}{2}\PY{p}{,} \PY{l+m+mi}{2}\PY{p}{,} \PY{l+m+mi}{6}\PY{p}{,} \PY{l+m+mi}{3}\PY{p}{,} \PY{n}{pi}\PY{o}{/}\PY{l+m+mi}{4}\PY{p}{)}
\PY{n}{parallelogram1}\PY{o}{.}\PY{n}{description}\PY{p}{(}\PY{p}{)}
\PY{n}{parallelogram1}\PY{o}{.}\PY{n}{print\PYZus{}superficy}\PY{p}{(}\PY{p}{)}
\PY{n+nb}{print}\PY{p}{(}\PY{l+s+sa}{f}\PY{l+s+s1}{\PYZsq{}}\PY{l+s+s1}{the superficy of the }\PY{l+s+si}{\PYZob{}}\PY{n}{parallelogram1}\PY{o}{.}\PY{n+nv+vm}{\PYZus{}\PYZus{}class\PYZus{}\PYZus{}}\PY{o}{.}\PY{n+nv+vm}{\PYZus{}\PYZus{}name\PYZus{}\PYZus{}}\PY{o}{.}\PY{n}{lower}\PY{p}{(}\PY{p}{)}\PY{l+s+si}{\PYZcb{}}\PY{l+s+s1}{ is: }\PY{l+s+si}{\PYZob{}}\PY{n}{parallelogram1}\PY{o}{.}\PY{n}{superficy}\PY{l+s+si}{\PYZcb{}}\PY{l+s+s1}{\PYZsq{}}\PY{p}{)}
\PY{n}{parallelogram1}\PY{o}{.}\PY{n}{draw\PYZus{}shape}\PY{p}{(}\PY{p}{)}
\end{Verbatim}
\end{tcolorbox}

    \begin{Verbatim}[commandchars=\\\{\}]
This is a parallelogram with the following properties:
Length:          6
Width:           3
Angle (in degrees):      45.0
The superficy of the parallelogram is: 12.727922061357857
the superficy of the parallelogram is: 12.727922061357857
    \end{Verbatim}

    \begin{center}
    \adjustimage{max size={0.9\linewidth}{0.9\paperheight}}{output_1_1.png}
    \end{center}
    { \hspace*{\fill} \\}
    
    In this example, we define a base class called \texttt{Shape} with an
\texttt{\_\_init\_\_} method that takes \texttt{x} and \texttt{y}
coordinates as parameters. We also define a \texttt{draw\_shape} method
that is empty and will be overridden in the subclass.\\
The \texttt{Parallelogram} class inherits from \texttt{Shape} and has
its own \texttt{\_\_init\_\_} method that takes additional parameters
for width, height, and angle. It also calculates the superficy of the
parallelogram. The \texttt{draw\_shape} method is overridden to provide
a specific implementation for drawing a parallelogram using Matplotlib.

Finally, we create an instance of the \texttt{Parallelogram} class and
call the \texttt{draw\_shape} method to visualize it.

    \begin{tcolorbox}[breakable, size=fbox, boxrule=1pt, pad at break*=1mm,colback=cellbackground, colframe=cellborder]
\prompt{In}{incolor}{3}{\boxspacing}
\begin{Verbatim}[commandchars=\\\{\}]
\PY{c+c1}{\PYZsh{} Define a class for rectangles that inherits from the parallelogram class, }
\PY{c+c1}{\PYZsh{} since a rectangle is a special case of a parallelogram where the angle is always 90 degrees (pi/2 radians).}
\PY{k}{class} \PY{n+nc}{Rectangle}\PY{p}{(}\PY{n}{Parallelogram}\PY{p}{)}\PY{p}{:}
    \PY{k}{def} \PY{n+nf+fm}{\PYZus{}\PYZus{}init\PYZus{}\PYZus{}}\PY{p}{(}\PY{n+nb+bp}{self}\PY{p}{,} \PY{n}{x0}\PY{p}{:}\PY{n+nb}{float}\PY{p}{,} \PY{n}{y0}\PY{p}{:}\PY{n+nb}{float}\PY{p}{,} \PY{n}{length}\PY{p}{:}\PY{n+nb}{float}\PY{p}{,} \PY{n}{width}\PY{p}{:}\PY{n+nb}{float}\PY{p}{)}\PY{p}{:}
        \PY{n+nb}{super}\PY{p}{(}\PY{p}{)}\PY{o}{.}\PY{n+nf+fm}{\PYZus{}\PYZus{}init\PYZus{}\PYZus{}}\PY{p}{(}\PY{n}{x0}\PY{p}{,} \PY{n}{y0}\PY{p}{,} \PY{n}{length}\PY{p}{,} \PY{n}{width}\PY{p}{,} \PY{n}{pi}\PY{o}{/}\PY{l+m+mi}{2}\PY{p}{)}   

    \PY{c+c1}{\PYZsh{} Define a method to print the description of the rectangle, including its properties: length and width }
    \PY{c+c1}{\PYZsh{} (we can omit the angle since it\PYZsq{}s always 90 degrees for a rectangle)}
    \PY{c+c1}{\PYZsh{} We can also call the description method of the parent class to reuse the code for printing length and width,}
    \PY{c+c1}{\PYZsh{} but we will override it here to provide a more specific description for rectangles.}
    \PY{k}{def} \PY{n+nf}{description}\PY{p}{(}\PY{n+nb+bp}{self}\PY{p}{)}\PY{p}{:} 
        \PY{n+nb}{print}\PY{p}{(}\PY{l+s+sa}{f}\PY{l+s+s2}{\PYZdq{}}\PY{l+s+s2}{This is a }\PY{l+s+si}{\PYZob{}}\PY{n+nb+bp}{self}\PY{o}{.}\PY{n+nv+vm}{\PYZus{}\PYZus{}class\PYZus{}\PYZus{}}\PY{o}{.}\PY{n+nv+vm}{\PYZus{}\PYZus{}name\PYZus{}\PYZus{}}\PY{o}{.}\PY{n}{lower}\PY{p}{(}\PY{p}{)}\PY{l+s+si}{\PYZcb{}}\PY{l+s+s2}{ with the following properties:}\PY{l+s+s2}{\PYZdq{}}\PY{p}{)}
        \PY{n+nb}{print}\PY{p}{(}\PY{l+s+s2}{\PYZdq{}}\PY{l+s+s2}{Length: }\PY{l+s+se}{\PYZbs{}t}\PY{l+s+s2}{\PYZdq{}}\PY{p}{,} \PY{n+nb+bp}{self}\PY{o}{.}\PY{n}{length}\PY{p}{)}        
        \PY{n+nb}{print}\PY{p}{(}\PY{l+s+s2}{\PYZdq{}}\PY{l+s+s2}{Width: }\PY{l+s+se}{\PYZbs{}t}\PY{l+s+se}{\PYZbs{}t}\PY{l+s+s2}{\PYZdq{}}\PY{p}{,} \PY{n+nb+bp}{self}\PY{o}{.}\PY{n}{width}\PY{p}{)}        

\PY{c+c1}{\PYZsh{} Define a class for squares that inherits from the Rectangle class, }
\PY{c+c1}{\PYZsh{} since a square is a special case of a rectangle where the length and width are equal.       }
\PY{k}{class} \PY{n+nc}{Square}\PY{p}{(}\PY{n}{Rectangle}\PY{p}{)}\PY{p}{:}
    \PY{k}{def} \PY{n+nf+fm}{\PYZus{}\PYZus{}init\PYZus{}\PYZus{}}\PY{p}{(}\PY{n+nb+bp}{self}\PY{p}{,} \PY{n}{x0}\PY{p}{:}\PY{n+nb}{float}\PY{p}{,} \PY{n}{y0}\PY{p}{:}\PY{n+nb}{float}\PY{p}{,} \PY{n}{side}\PY{p}{:}\PY{n+nb}{float}\PY{p}{)}\PY{p}{:}
        \PY{n+nb}{super}\PY{p}{(}\PY{p}{)}\PY{o}{.}\PY{n+nf+fm}{\PYZus{}\PYZus{}init\PYZus{}\PYZus{}}\PY{p}{(}\PY{n}{x0}\PY{p}{,} \PY{n}{y0}\PY{p}{,} \PY{n}{side}\PY{p}{,} \PY{n}{side}\PY{p}{)}
        \PY{n+nb+bp}{self}\PY{o}{.}\PY{n}{side}\PY{o}{=}\PY{n}{side}  \PY{c+c1}{\PYZsh{} Initialize the side length of the square, which is equal to both the length and width of the rectangle}
    
    \PY{c+c1}{\PYZsh{} Define a method to print the description of the square, including its properties: side length}
    \PY{c+c1}{\PYZsh{} We can also call the description method of the parent class to reuse the code for printing length and width,}
    \PY{c+c1}{\PYZsh{} but we will override it here to provide a more specific description for squares.}
    \PY{k}{def} \PY{n+nf}{description}\PY{p}{(}\PY{n+nb+bp}{self}\PY{p}{)}\PY{p}{:} 
        \PY{n+nb}{print}\PY{p}{(}\PY{l+s+sa}{f}\PY{l+s+s2}{\PYZdq{}}\PY{l+s+s2}{This is a }\PY{l+s+si}{\PYZob{}}\PY{n+nb+bp}{self}\PY{o}{.}\PY{n+nv+vm}{\PYZus{}\PYZus{}class\PYZus{}\PYZus{}}\PY{o}{.}\PY{n+nv+vm}{\PYZus{}\PYZus{}name\PYZus{}\PYZus{}}\PY{o}{.}\PY{n}{lower}\PY{p}{(}\PY{p}{)}\PY{l+s+si}{\PYZcb{}}\PY{l+s+s2}{ with the following properties:}\PY{l+s+s2}{\PYZdq{}}\PY{p}{)}
        \PY{n+nb}{print}\PY{p}{(}\PY{l+s+s2}{\PYZdq{}}\PY{l+s+s2}{Side: }\PY{l+s+se}{\PYZbs{}t}\PY{l+s+s2}{\PYZdq{}}\PY{p}{,} \PY{n+nb+bp}{self}\PY{o}{.}\PY{n}{side}\PY{p}{)}        


\PY{c+c1}{\PYZsh{} Create an instance of the parallelogram class and call its methods to demonstrate functionality}
\PY{c+c1}{\PYZsh{} populating the parameters with some example values (x0=2, y0=2, length=9, width=4, angle=pi/6 : 30 degrees)}
\PY{n}{parallelogram1} \PY{o}{=} \PY{n}{Parallelogram}\PY{p}{(}\PY{l+m+mi}{2}\PY{p}{,} \PY{l+m+mi}{2}\PY{p}{,} \PY{l+m+mi}{9}\PY{p}{,} \PY{l+m+mi}{4}\PY{p}{,} \PY{n}{pi}\PY{o}{/}\PY{l+m+mi}{6}\PY{p}{)}
\PY{n}{parallelogram1}\PY{o}{.}\PY{n}{description}\PY{p}{(}\PY{p}{)}
\PY{n+nb}{print}\PY{p}{(}\PY{l+s+s1}{\PYZsq{}}\PY{l+s+s1}{the superficy of the parallelogram is:}\PY{l+s+s1}{\PYZsq{}}\PY{p}{,} \PY{n}{parallelogram1}\PY{o}{.}\PY{n}{superficy}\PY{p}{)}
\PY{n}{parallelogram1}\PY{o}{.}\PY{n}{draw\PYZus{}shape}\PY{p}{(}\PY{p}{)}

\PY{c+c1}{\PYZsh{} Create an instance of the Rectangle class and call its methods to demonstrate functionality}
\PY{n}{rectangle1} \PY{o}{=} \PY{n}{Rectangle}\PY{p}{(}\PY{l+m+mi}{1}\PY{p}{,} \PY{l+m+mi}{1}\PY{p}{,} \PY{l+m+mi}{10}\PY{p}{,} \PY{l+m+mi}{4}\PY{p}{)}
\PY{n}{rectangle1}\PY{o}{.}\PY{n}{description}\PY{p}{(}\PY{p}{)}
\PY{n}{rectangle1}\PY{o}{.}\PY{n}{print\PYZus{}superficy}\PY{p}{(}\PY{p}{)}
\PY{n+nb}{print}\PY{p}{(}\PY{l+s+s1}{\PYZsq{}}\PY{l+s+s1}{the superficy of the rectangle is:}\PY{l+s+s1}{\PYZsq{}}\PY{p}{,} \PY{n}{rectangle1}\PY{o}{.}\PY{n}{superficy}\PY{p}{)}
\PY{n}{rectangle1}\PY{o}{.}\PY{n}{draw\PYZus{}shape}\PY{p}{(}\PY{p}{)}

\PY{c+c1}{\PYZsh{} Create an instance of the Square class and call its methods to demonstrate functionality}
\PY{n}{square1} \PY{o}{=} \PY{n}{Square}\PY{p}{(}\PY{l+m+mi}{3}\PY{p}{,} \PY{l+m+mi}{3}\PY{p}{,} \PY{l+m+mi}{3}\PY{p}{)}
\PY{n}{square1}\PY{o}{.}\PY{n}{description}\PY{p}{(}\PY{p}{)}
\PY{n}{square1}\PY{o}{.}\PY{n}{print\PYZus{}superficy}\PY{p}{(}\PY{p}{)}
\PY{n+nb}{print}\PY{p}{(}\PY{l+s+s1}{\PYZsq{}}\PY{l+s+s1}{the superficy of the square is:}\PY{l+s+s1}{\PYZsq{}}\PY{p}{,} \PY{n}{square1}\PY{o}{.}\PY{n}{superficy}\PY{p}{)}
\PY{n}{square1}\PY{o}{.}\PY{n}{draw\PYZus{}shape}\PY{p}{(}\PY{p}{)}
\end{Verbatim}
\end{tcolorbox}

    \begin{Verbatim}[commandchars=\\\{\}]
This is a parallelogram with the following properties:
Length:          9
Width:           4
Angle (in degrees):      29.999999999999996
the superficy of the parallelogram is: 17.999999999999996
    \end{Verbatim}

    \begin{center}
    \adjustimage{max size={0.9\linewidth}{0.9\paperheight}}{output_3_1.png}
    \end{center}
    { \hspace*{\fill} \\}
    
    \begin{Verbatim}[commandchars=\\\{\}]
This is a rectangle with the following properties:
Length:          10
Width:           4
The superficy of the rectangle is: 40.0
the superficy of the rectangle is: 40.0
    \end{Verbatim}

    \begin{center}
    \adjustimage{max size={0.9\linewidth}{0.9\paperheight}}{output_3_3.png}
    \end{center}
    { \hspace*{\fill} \\}
    
    \begin{Verbatim}[commandchars=\\\{\}]
This is a square with the following properties:
Side:    3
The superficy of the square is: 9.0
the superficy of the square is: 9.0
    \end{Verbatim}

    \begin{center}
    \adjustimage{max size={0.9\linewidth}{0.9\paperheight}}{output_3_5.png}
    \end{center}
    { \hspace*{\fill} \\}
    
    The description of the code is as follows:

\begin{enumerate}
\def\labelenumi{\arabic{enumi}.}
\item
  We import the necessary libraries: \texttt{matplotlib.pyplot} for
  plotting and \texttt{numpy} for numerical operations.
\item
  We define the \texttt{Shape} class with an \texttt{\_\_init\_\_}
  method that initializes the \texttt{x} and \texttt{y} coordinates and
  a placeholder method \texttt{draw\_shape}.
\item
  We define the \texttt{Parallelogram} class that inherits from
  \texttt{Shape}. It has its own \texttt{\_\_init\_\_} method that
  initializes additional attributes for width, height, and angle. It
  also calculates the superficy of the parallelogram.
\item
  The \texttt{draw\_shape} method in the \texttt{Parallelogram} class
  calculates the vertices of the parallelogram based on the given
  parameters and uses Matplotlib to draw it.
\item
  The \texttt{description} method in the \texttt{Parallelogram} class
  provides a textual description of the parallelogram, including its
  dimensions and area.
\item
  we create an instance of the \texttt{Parallelogram} class and call the
  \texttt{draw\_shape} method to visualize the parallelogram.
\item
  We define the \texttt{Rectangle} class that inherits from
  \texttt{Parallelogram}. It has its own \texttt{\_\_init\_\_} method
  that initializes the length and width of the rectangle, and it
  overrides the \texttt{description} method to provide a more specific
  description for rectangles.
\item
  We define the \texttt{Square} class that inherits from
  \texttt{Rectangle}. It has its own \texttt{\_\_init\_\_} method that
  initializes the side length of the square, and it overrides the
  \texttt{description} method to provide a more specific description for
  squares.
\item
  Finally, we create instances of the \texttt{Rectangle} and
  \texttt{Square} classes and call their \texttt{description} methods to
  see the specific details of each shape.
\item
  The output of the code will include the visualization of the
  parallelogram and the descriptions of the rectangle and square,
  demonstrating inheritance and polymorphism in Python OOP.
\item
  polymorphism is demonstrated through the \texttt{description} method,
  which is overridden in both the \texttt{Rectangle} and \texttt{Square}
  classes to provide specific details about each shape, while still
  being called in a similar way. it is also demonstrated throught the
  \texttt{calculate\_area} method, which is defined in the
  \texttt{Parallelogram} class and inherited by both the
  \texttt{Rectangle} and \texttt{Square} classes, allowing us to
  calculate the area of each shape using the same method.
\end{enumerate}


    % Add a bibliography block to the postdoc
    
    
    
\end{document}
